\documentclass[a4paper, 12pt, french]{article}

% ========================
% IMPORTATION DES PACKAGES
% ========================

\usepackage[T1]{fontenc}
\usepackage[utf8]{inputenc}
\usepackage[margin=2cm]{geometry}
\usepackage{lmodern}
\usepackage{theorem}
\usepackage{amsfonts, amsmath, amssymb}
\usepackage{babel}

\pagestyle{headings}

\renewcommand{\familydefault}{\sfdefault}

\theoremstyle{break}
\theoremheaderfont{\scshape}
\theorembodyfont{\upshape}

\newtheorem{enon}{Énoncé}[section]
\newtheorem{preuve}{Preuve}[section]

\title{MÉSURE ET INTÉGRATION}
\author{}
\date{\today}

\begin{document}
\maketitle

\begin{center}\section{Théorie des ensembles}\end{center}
\begin{enon}
       Soit $(A_{n})_{n \in \mathbb{N}} \subseteq P(X)$.
       \begin{enumerate}
               \item Montrer que  $\displaystyle\bigcup_{n \in \mathbb{N}} A_{n} =  \bigsqcup_{n \in \mathbb{N}} B_{n}$.
               \item Est ce qu'on peut écrire  $\displaystyle\bigcup_{n \in \mathbb{N} } A_{n}$ comme une suite croissante ? respectivement comme une suite décroissante ?
       \end{enumerate}
\end{enon}
\begin{preuve}
      \begin{enumerate}
              \item Soit $(A_{n})_{n \in \mathbb{N}} \subseteq P(X)$.\\ Posons $ B_{1}= A_{1}, \; B_{n} = A_{n} \backslash \bigcup_{n=1}^{n-1} A_{n}$ pour $n \geq 2$. On peut voir facilement que les $B_{n}$ sont deux à deux disjoints.\\
              D'abord montrons que $\displaystyle\bigcup_{n \in \mathbb{N}} A_{n} \subseteq \bigsqcup_{n \in \mathbb{N}} B_{n}$.\\
              Soit $ x \in \displaystyle\bigcup_{n \in \mathbb{N}} A_{n}$, alors il existe $n_{0}$ tel que $x \in A_{n_{0}}$. Alors $x \not\in \displaystyle\bigcup_{k<n_{0}} A_{k}$. Donc $x \in A_{n_{0}} \backslash\displaystyle\bigcup_{n \in \mathbb{N}} A_{n} = B_{n_{0}}$.
              Ainsi $x \in \displaystyle\bigsqcup_{n \in \mathbb{N}} B_{n}$.\\
              Montrons maintenant que $\displaystyle\bigsqcup_{n \in \mathbb{N}} B_{n} \subseteq  \displaystyle\bigcup_{n \in \mathbb{N}} A_{n}$.\\
              Soit $x \in \bigsqcup_{n \in \mathbb{N}} B_{n}$, alors il existe $n_{0} \in \mathbb{N}$ tel que $x \in B_{n_{0}} = A_{n_{0}} \backslash \bigcup_{n_{0}=1}^{n-1}A_{n_{0}}$. Ainsi $x \in A_{n_{0}}$. Donc $x \in \bigcup_{n \in \mathbb{N}}$. D'où $\displaystyle\bigsqcup_{n \in \mathbb{N}} B_{n} \subseteq  \displaystyle\bigcup_{n \in \mathbb{N}} A_{n}$.
              On conclut alors que  $\displaystyle\bigcup_{n \in \mathbb{N}} A_{n} =  \bigsqcup_{n \in \mathbb{N}} B_{n}$.
              
              \item $(A_{n})_{n \in \mathbb{N}} \subseteq P(X)$.\\ Oui, on peut écrire  $\displaystyle\bigcup_{n \in \mathbb{N} } A_{n}$ comme une suite croissante.
              Si on pose $C_{k} = \displaystyle\bigcup_{n \in \mathbb{N}} A_{n}$. $C_{k} \subseteq C_{k+1}$. Ainsi $\displaystyle\bigcup_{n \in \mathbb{N}} C_{n} = \bigcup_{n \in \mathbb{N}} A_{n}$. De même pour l'indice $n+1$. Ainsi elle est démontrée.
              En général, on ne peut pas écrire $\displaystyle\bigcup_{n \in \mathbb{N}} A_{n}$ comme réunion d'une suite décroissante.\\ Une suite décroissante vérifie $D_{n+1} \subseteq D_{n}$. Dans ce cas $\displaystyle\bigcup_{n \in \mathbb{N}} D_{n} = D_{0}$. Donc la réunion d'une suite décroissante est déjà fixé dès le premier terme.
      \end{enumerate}
\end{preuve}

\begin{enon}[Injectivite, surjectivité et bijectivité]
	Soient $X$ et $Y$ deux ensemble non vide et $f : X \rightarrow Y$ une application. Montrer que
	\begin{itemize}
		\item[a)] $f$ est injective si, et seulement si pour tout $A \in P(X)$, on a $f^{-1}(f(A)) = A$.
		\item[b)] $f$ est surjective si, et seulement si pour tout $A \in P(Y)$, on a $f(f^{-1}(B)) = B$.
		\item[a)] $f$ est bijective si, et seulement si pour tout $A \in P(X)$, on a $f(X \backslash A)=Y \backslash f(A)$.
	\end{itemize}
\end{enon}

\begin{preuve}
	Soient $X$ et $Y$ deux ensemble non vide et $f : X \rightarrow Y$ une application.
	\begin{itemize}
		\item[a)] Montrons que $f$ est injective si, et seulement si pour tout $A \in P(X)$, on a $f^{-1}(f(A)) = A$.\\
		Supposons que $f$ est injective. Pour tout $A \in P(X)$, nous allons montrer que $f^{-1}(f(A)) = A$.\\
		D'abord montrons que $A \subseteq f^{-1}(f(A))$, soit $x \in A$. Par définition $f(x) \in f(A)$, or $x \in f^{-1}(f(A))$. D'où $A \subseteq f^{-1}(f(A))$.\\ Ensuite montrons que $f^{-1}(f(A)) \subseteq A$, soit $x \in f^{-1}f(A)$. Par défintion $f(x) \in f(A)$, il existe $x^{'} \in A$ tel que $f(x)=f(x^{'})$ cela implique que $x = x^{'}$ puisque $f$ est injective. Alors $x \in A$. D'où $f^{-1}(f(A)) \subseteq A$.\\ Reciproquement, soit $A \in P(X)$ et on a $f^{-1}(f(A)) = A$. Soient $x,y \in A$ tel que $f(x)=f(y)$. Considerons l'ensemble $A=\{x\}$. Par hypothèse, on a $f^{-1}f(\{x\})=\{x\}$, alors $f(\{x\}) \in f(A)$. On a $f(x)=f(y)$, donc $f(y) \in f(A)$. Par définition de la fonction reciproque, on a $y \in f^{-1}f(A)$ alors $y \in \{x\}$. Donc $x=y$. D'où l'injectivité de $f$. \marginpar{$\square$} 
	\end{itemize}
\end{preuve}

\begin{center}\section{Tribus, clan et classe monotone}\end{center}
\begin{enon}
      Soit $X$ un ensemble non dénombrable. On pose $$\tau \{ A \in P(X) | A \mbox{ est dénombrable ou } X \backslash A \mbox{ est dénombrable}\}.$$
      \begin{enumerate}
              \item Montrer que $\tau$ est un tribu sur $X$.
              \item Comparer le tribu engendré par $\{x\}$ et $\tau$.
      \end{enumerate}
\end{enon}

\begin{preuve}
   Soit $X$ un ensemble non dénombrable. On pose $$\tau := \{ A \in P(X) | A \mbox{ est dénombrable ou } X \backslash A \mbox{ est dénombrable}\}.$$
   \begin{enumerate}
          \item Montrons que $\tau$ est un tribu sur $X$.
          \begin{itemize}
                     \item On sait que $\emptyset$ et $X \in P(X)$. $\emptyset$ est dénombrable car elle est finie. $X \backslash X = \emptyset$, alors $X\backslash X$ est aussi dénombrable. D'où $\emptyset, X \in \tau$.
                     \item Soit $A \in \tau$, alors $A \in P(X)$. Par définition de $\tau$, $X\backslash A$ est dénombrable. Donc $X \backslash A \in \tau$. Elle est alors stable par passage au complémentaire.
                     \item Soit $(A_{n})_{n \in \mathbb{N}} \in \tau$. Deux cas s'offrent à nous 
                     \begin{itemize}
                             \item[Cas 1: ] Si $A_{n}$ est dénombrable. \\
                             Pour tout $n \in \mathbb{N}$, $\bigcup_{n} A_{n}$ est dénombrable car union d'ensemble dénombrable est dénombrable. Donc $\bigcup_{n} A_{n} \in \tau$.
                             \item[Cas 2: ] Si $X\backslash A_{n}$ est dénombrable. \\
                             Pour tout $n \in \mathbb{N}$, $\bigcup_{n} X\backslash A_{n}$ est dénombrable. Donc $\bigcup_{n} X\backslash A_{n} \in \tau$.
                     \end{itemize}
                     Dans les deux cas on a montré qu'elles sont stable par réunion dénombrable.
          \end{itemize}
          On conclut alors que $\tau$ est bien un tribu sur $X$.
   \end{enumerate}
\end{preuve}
\end{document}
